% Section II - Related work. AI DRAFT (Claude, 2026-09-24): Kunal and advisor to edit.
% Budget: ~380 total (option B) or ~300 (option A); see word_budget_MLESP.md, advisor Q2. Per blueprint Sec 9, cut here and in Section I first if long.
% REVIEW 2026-09-27 (C3): wen2020 and young2025 added as design precedents; Yagis random-label clause cut to offset.
% II-A describes corpus studies from their full texts (read during screening). II-B: kaur2025 and zhang2025 are
% described from their ABSTRACTS only (Europe PMC, checked 2026-09-25; full texts are Elsevier, not opened).
% II-C was checked by Claude against the Europe PMC open-access full texts of lee2023 (PMC10547830),
% rosenblatt2024 (PMC10901797) and yagis2021 (PMC8604922) on 2026-09-24. Blueprint Sec 12 still requires Kunal
% to open each original himself; their bibliography entries stay UNVERIFIED until he does.
% 2026-09-25: II-C corrected after KC's own reference check (every correction re-verified by Claude against the
% same full texts): Rosenblatt's leaky covariate regression DEFLATED performance; Yagis's 66->97%
% is the VGG16/OASIS-200 case; Lee's ranges are per model.

\subsection{Deep learning on EEG connectivity for MDD}
In this literature a connectivity measure turns each EEG segment into a channel-by-channel matrix, and a deep network classifies the matrices. Most measures are undirected phase or amplitude couplings (PLV, PLI, coherence, correlation and mutual information) \cite{li2020,dang2020,zhang2023,yu2022,liu2024}; a smaller group uses directed measures from multivariate autoregressive models \cite{saeedi2021,khan2021}. Some stack one matrix per frequency band into a multilayer network \cite{dang2020,sun2022}. Convolutional networks treat the matrix as an image, sometimes with a recurrent branch \cite{saeedi2021,zhang2023,li2024}; graph neural networks treat channels as nodes and connectivity as edge weights \cite{sun2022,sun2024,liu2024}, and recent work learns the graph during training \cite{luo2023,xia2023} or moves it to cortical source space \cite{guo2025}. A few public datasets with tens of participants per group recur (Table~\ref{tab:corpus}). Reported accuracies commonly exceed 95\%, yet these studies seldom state when the pipeline's statistics were computed relative to the train/test split.

\subsection{Existing reviews of validation practice}
% REVIEW 2026-11-07 (C10 space pass, blueprint Sec 9: cut from II before III or V): condensed to free
% room for the reviewer-requested analyses in Sections IV, V and VI. No citation or claim dropped.
Recent reviews compare validation schemes in EEG depression studies; one found none that validated across datasets \cite{lin2024}, another found near-perfect accuracy most often where samples were small and validation subject-dependent \cite{lunaguevara2026}, and a third found that the lack of standardized evaluation protocols limits comparison \cite{kaur2025}. A meta-analysis of 176 machine-learning studies of MDD found reported accuracy negatively correlated with sample size \cite{zhang2025}---the confound this audit addresses through sample-size stratification.

The distinction we draw is one of scope. These reviews ask how a study separated training from test data, and whether its validation was subject-dependent or subject-independent, and one of them also discusses feature selection \cite{lunaguevara2026}. None that we found asks whether the steps that build the connectivity matrix itself---estimation, thresholding and normalization---were fit before that separation. A study can use leave-one-subject-out validation and still leak through a threshold chosen on the pooled data; a review that records only the validation scheme would score it as sound.
% REVIEW 2026-11-07: the \PENDING on this sentence is resolved conservatively. Rather than assert that no
% review covers L5, the claim is narrowed to L2-L4 and lunaguevara2026 is credited for feature selection, so
% the sentence holds whether or not that review's full text covers L5. Nothing here needed new verification.

\subsection{Data leakage in neuroimaging machine learning}
Leakage is well documented in neuroimaging. Across four resting-state fMRI connectome datasets, selecting features on the full dataset and repeating subjects across the split inflated performance most \cite{rosenblatt2024}. In MRI classification, splitting by slice rather than subject raised a VGG16 model's accuracy on one dataset from 66\% to about 97\% \cite{yagis2021}. In resting-state EEG from patients with post-traumatic stress disorder, splitting cropped trials rather than participants inflated balanced accuracy by about 5--15 percentage points for one convolutional network and 15--20 for EEGNet \cite{lee2023}. Segment-based rather than subject-based holdout strongly overestimated accuracy on unseen subjects, and most translational EEG studies surveyed used it \cite{brookshire2024}; in EEG-based schizophrenia detection, about 65\% of studies split by epoch or selected features before partitioning, inflating accuracy by up to about 30\% \cite{racz2026}. None examined how connectivity features are built. The connectome work is the closest analogue, deriving features from pairwise connectivity; the two forms it found most damaging correspond to L5 and L1 here. Rating leakage from reported methods also has precedents: one survey labeled each Alzheimer's study's leakage status from its text \cite{wen2020}, and a scoping review stratified 44 such studies into three risk tiers, finding reported accuracy highest where risk was highest \cite{young2025}. We adapt that approach to the steps that build a connectivity matrix.
